\newpage
\section{Nicht-kohärenter 2-FSK-Empfänger}

% \begin{itemize}[noitemsep]
%     \item Kurze Beschreibung der Aufgabenstellung
%     \item Blockdiagramm des implementierten nicht-kohärenten Empfängers
%     \item Screenshot des GNU Radio Companion Aufbaus
%     \begin{itemize}[noitemsep]
%         \item Parameter der Bandpassfilter (warum wurden diese so gewählt?)
%     \end{itemize}
%     \item Analyse des empfangenen Signals im Zeitbereich
%     \begin{itemize}[noitemsep]
%         \item Grafische Darstellung
%         \item Interpretation des Ergebnisses 
%     \end{itemize}
%     \item Analyse des empfangenen Signals im Frequenzbereich
%     \begin{itemize}[noitemsep]
%         \item Grafische Darstellung
%         \item interpretation des ergebnisses
%     \end{itemize}
% \end{itemize}

Es ist ein nicht-kohärenter 2-FSK Empfänger zu implementieren, der das Signal aus \autoref{sec:task2} demoduliert. Dazu wird der Aufbau wie in \autoref{fig:noncoherent_fsk_receiver} gezeigt in GRC implementiert. Anschließend wird das empfangene Signal im Zeit- und Frequenzbereich untersucht.

\subsection{Blockschaltbild eines nicht-kohärenten 2-FSK Empfängers}

\begin{figure}[h]
    \centering
    \begin{circuitikz}[>=latex', transform shape, scale=0.9]
        \draw (0,0) node[bareRXantenna](ant){} |- ++(1,-0.5) coordinate (aout) node[circ]{} node[right]{$r[k]\in\mathbb{C}$};
        \node[above] at (ant.top){USRP};
        \draw (aout) |- ++(1,1) to[bandpass,>,l^=$f_0+\frac{f_\text{dev}}{2}$]
            ++(2,0) to[twoport,>,align=center,l={Envelope\\Detector}, t={$(\cdot)^2$}]
            ++(2,0) -- ++(1.5,0) coordinate (sum1) node[anchor=south east]{$r_1[k]\in\mathbb{R}$};
        \draw (aout) |- ++(1,-1) to[bandpass,>,l_=$f_0-\frac{f_\text{dev}}{2}$]
            ++(2,0) to[twoport,>,align=center,t={$(\cdot)^2$}]
            ++(2,0) -- ++(1.5,0) coordinate (sum2) node[anchor=north east]{$r_0[k]\in\mathbb{R}$};
        \node[adder, scale=0.7](add) at ($(sum1)!0.5!(sum2)$) {};
        \draw[->] (sum1) -- (add.north) node[pos=0.5, left] {$+$};
        \draw[->] (sum2) -- (add.south) node[pos=0.5, left] {$-$};
        \draw[->] (add.east) to[allornothing, align=center, l={Threshold\\Detector}] ++(2,0) node[right, text width=5em] {Bitstream Output};
    \end{circuitikz}
    \caption{Blockschaltbild eines nicht-kohärenten FSK Empfängers}\label{fig:noncoherent_fsk_receiver}
\end{figure}

Das Empfangene Signal kann nun nicht mehr wie in \autoref{sec:task2} symmetrisch um \SI{0}{\hertz} gefiltert werden. Je Symbol muss der jeweilige Frequenzanteil mit einem Bandpassfilter herausgefiltert werden, damit er anschließend mit einem Hüllkurvendemodulator detektiert werden kann.

Die Mittelfreqenz der Bandpässe wird so gewählt, dass die Frequenzen des \texttt{'1'} Symbols bzw. des \texttt{'0'} Symbols durchgelassen werden. Eine Aufzeichnung des Spektrum des empfangenen Basisbandsignals ergab, dass die frequenzanteile bei $\pm \frac{f_\mathrm{dev}}{2}$ liegen. Eine Bandbreite von ca. $\SI{10}{\kilo\hertz}$ ist ausreichend, um die Signale durchzulassen.

\subsection{Aufbau in GRC}

\begin{figure}[h]
    \centering
    \includegraphics[width=\textwidth]{\noncoherentPath{3_grc.pdf}}
    \caption{GNU Radio Companion Aufbau für Task 3}
\end{figure}

Der Hüllkurvendetektor wird durch den \textit{Complex to Mag Squared} Block realisiert. Ob ein Symbol als \texttt{'1'} oder \texttt{'0'} erkannt wird, wird durch den \textit{Binary Slicer} Block entschieden.

\subsection{Analyse im Zeitbereich}

\begin{figure}[h]
    \centering
    \includegraphics[width=\textwidth]{\noncoherentPath{RecievedBitStream_Annotated.png}}
    \caption{Empfangener Bitstream am Ausgang des Demodulators}\label{fig:recieved_bitstream}
\end{figure}

Im Zeitbereich ist nun nicht mehr das empfangene Symbol $\in\{-1,1\}$ sichtbar, sondern der tatsächliche Bitstrom $\in\{0,1\}$. Gesendet wurde das Datenwort \texttt{0xB3} $\equiv$ \texttt{0b1011\,0011}. Es wurde Festgestellt, dass die zweige für das \texttt{-1}-Symbol und das \texttt{1}-Symbol vertauscht sind, daher sind in \autoref{fig:recieved_bitstream} die Pegel invertiert dargestellt. Trotzdem ist zu erkennen, dass das Datenwort korrekt empfangen wurde: Das Abgelesene Datenwort ist \texttt{0b0100\,1100}. Invertiert ergibt das \texttt{0b1011\,0011} $\equiv$ \texttt{0xB3}, was dem gesendeten Datenwort entspricht.

\subsection{Analyse im Frequenzbereich}

Ohne Filterung ist das gesamte Spektrum wie in \autoref{fig:baseband_no_bpf} dargestellt sichtbar. Für den Eingestellten Kanal (Kanal 1) treten die gewünschten Frequenzkomponenten um \SI{0}{\hertz} auf, jedoch ist auch die aktivität der Nachbarkanäle sichtbar, welche das Signal stören können.

\begin{figure}[h]
    \centering
    \includegraphics[width=\textwidth]{\noncoherentPath{Baseband.png}}
    \caption{Empfangenes Signal ohne Filterung}\label{fig:baseband_no_bpf}
\end{figure}

Mit den Bandpassfiltern ist in \autoref{fig:baseband_bpf} zu erkennen, dass nurnoch die gewünschten Frequenzanteile $\pm \frac{f_\mathrm{dev}}{2}$ durchgelassen werden. Der Hüllkurvendetektor kann somit die Amplitude der jeweiligen Frequenzanteile bestimmen, um die Symbole zu unterscheiden.

\begin{figure}[h]
    \centering
    \includegraphics[width=\textwidth]{\noncoherentPath{Baseband_BPF.png}}
    \caption{Empfangenes Signal mit Bandpassfilterung}\label{fig:baseband_bpf}
\end{figure}